\exnumber{00}
\extitle{UX II - Ex 1 - Guide d'entretien}
\exfiles{ex00/PDF, nomme Name\_FirstName\_DeliverableName\_Date\_VersionNumber}
\exauthorize{Those named in the exercise}
\makeheaderfiles

\textbf{Contexte :}

Welcome to the 2nd quest.

Your mission, should you choose to accept it, is to focus on going on holidays with friends. As you may have noticed, experiencing weekends or other kinds of holidays is super nice, but may lead to some moments of irritation.

So this project's goal is to make holidays more comfortable fo everyone. But first, as in all design processes, you need to understand the real problem behind stressful moments!

\textbf{Instructions :}

You need to create an interview guide to get information from your user.

And remember, you are not your user!

\textbf{Livrable :}

\begin{itemize}
\item A PDF document with an interview guide
\item Déposez le PDF sur Github, nommé ainsi : « Name\_FirstName\_DeliverableName\_Date\_VersionNumber ». Exemple : « Doe\_John\_MidFiPrototype\_05242024\_V1 ».
\item Ajoutez un titre dans le document.
\item Soignez l'aspect global du document. Il doit être clair, simple et lisible. Vous pouvez regarder canva.com, sans surcharger le dessin de détails.
\end{itemize}

\textbf{Pistes :}
Voici les étapes recommandées :

1. Define your objective, what you are looking for when interviewing people. Usually, this has to do with your user's habits, and the most painful problems they have regarding this topic.
2. Write down 1 to 3 hypotheses you have regarding this topic, that needs to be either confirmed or disproved.
3. Choisissez 3 à 5 sujets sur lesquels il vous faut de l'information.
4. For each topic, write 1 to 3 open-ended questions. Remember, your user cannot answer with "Yes" or "No!" Ex: How, Why, What, When, Tell me about the last time you did... etc. If you really need to introduce a new topic by a close-ended question, ask why right after.
5. Write some screening questions to have a few pieces of demographics (age, job, location, etc.)
6. Structure your questions in a funnel : start with broad questions to focus on specific topics or situations. Remember to stay neutral in your tone to let the user tell their own story. Try to focus on 10 main questions.
7. Write a greeting introduction, describing briefly (1) who you are and (2) what you are looking for. No need to say too much, but enough to provide a framework and the purpose of the interview.
8. Write a conclusion to thank your user for their time, and prepare the next steps. Ex: ask to be put in touch with someone they know and who could be interested in this topic, etc.

\textbf{Ressources :}

\begin{itemize}
\item Respondent - Tool to find users \url{https://www.respondent.io}
\item How to conduct user interviews \url{https://uxdesign.cc/how-to-conduct-user-interviews-fe4b8c34b0b7?sk=9625632a522a9309a91297191e668565}
\item How to design better products through user interviews \url{https://uxdesign.cc/how-to-design-better-products-through-user-interviews-4c5142bb1fc4}
\item Asking the right questions \url{https://uxdesign.cc/asking-the-right-questions-on-user-research-interviews-and-testing-427261742a67}
\item Podcast recommendation: How to Ask Great UX Research Questions \url{https://open.spotify.com/episode/5UlvTKyVObeCDxkc8DAakU?si=fjMEHCdjRXWkYhwKhxIopQ}
\end{itemize}

\textbf{Citation :}

"To ask open-ended questions is the best approach, but it's easy to get into the weeds in data analysis when every answer is a paragraph or two of prose. Users quickly tire of answering many open-ended questions, which usually require a lot of typing and explanation." Norman Nielsen Group

